% This manuscript is untagged; claim archival PDF/A-1b, not accessible PDF/A-1a.
% pdfx processes declared options in order, with a-1b after the class's a-1a.
\PassOptionsToPackage{a-1b}{pdfx}
\documentclass{ceurart}
\usepackage{longtable}
\usepackage{bookmark}
\providecommand{\tightlist}{\setlength{\itemsep}{0pt}\setlength{\parskip}{0pt}}
\providecommand{\pandocbounded}[1]{#1}
\begin{document}
\copyrightyear{2027}
\conference{IUI Workshops 2027: Joint Proceedings of the ACM Conference on Intelligent User Interfaces Workshops; HIMAS, February 2027, Helsinki, Finland}
\title{REAL: Human-Inspectable Intermediate Representations for AI-Assisted Video Editing}
\author{Anonymous authors}
\begin{abstract}
An incorrect video edit can arise from a mistaken media description, an
unsuitable selection or a timing error. An editable project alone does
not identify which decision needs repair. REAL organizes AI-assisted
editing around four human-inspectable representations: a narrative
notebook, media catalogue, authoritative storyboard and editable
project. Word anchors, decision ownership and explicit revision commands
make choices addressable and record corrections for rebuilding and
export. We examine the technical conditions for this interaction
contract. Across 30 development and 100 additional synthetic projects,
the final frozen implementation reproduces every reference placement
under declared assembly semantics. Re-anchoring one overlay in each
development project preserves the other event identities and boundaries;
corrections also survive rebuilding under two artificial time scales on
the same 100-project cohort. In a retrospective comparison, two models
supplied with the same plans and format notes also reproduce all 30
projects exactly, locating REAL's contribution in addressable decisions
and their correction. Omitting format notes exposes format errors
despite close temporal scores. A narrated reconstruction shows a
complementary limitation: correct asset selection does not resolve pause
and caption conventions. Effects on creators' understanding, effort and
agency remain untested. REAL contributes a working inspectable workflow,
reproducible conformance evidence and design implications for making
multimodal decisions locally correctable.

\end{abstract}
\begin{keywords}
AI-assisted video editing \sep human-AI collaboration \sep intermediate representations \sep multimodal systems \sep local correction
\end{keywords}
\maketitle
\section{Introduction}\label{introduction}

A creator assembling a narrated video coordinates what the narration
means, what the footage contains, which source interval supports each
point, and when images, titles and sound should appear. A misplaced
insert may reflect a wrong description, an unsuitable selection or a
correct choice translated into the wrong time. Repair becomes difficult
when these decisions are collapsed into one generated deliverable.

Consider a violin lesson in which a fingering diagram should appear as
the player explains an octave shift. If the diagram arrives late, the
creator needs to move it to the relevant phrase while keeping the
selected footage, caption style and other placements intact.

Interactive editing systems already connect transcripts to footage
\citep{truong2016quickcut, huber2019bscript}, describe media in language
\citep{wang2024lave}, and expose persistent production artifacts
\citep{yan2026crayotter}. REAL builds on this direction with an
augmentation-oriented stance \citep{brynjolfsson2022turing}: the creator
remains the author, and intermediate choices should be available for
correction. Its systems question is: \textbf{How can a creator revise an
intermediate editing decision while the workflow preserves that
correction through timing and export?}

REAL treats editing as compilation through four readable artifacts. A
notebook expresses narrative intent; a catalogue records media evidence
and suggested uses; a storyboard owns editorial choices; a compiler
translates them into an editable Final Cut Pro project. Models interpret
media and requests, while scripts implement binding rules, timing,
revisions and validation. A stable word anchor can retain a placement's
identity when recorded speech replaces provisional timing; creator
ownership marks choices that a rebuild must retain.

Readable plans must survive export, while correctly exported footage may
still fail to communicate the idea. We distinguish structural
conformance, media selection and creator outcomes.

We contribute:

\begin{enumerate}
\def\labelenumi{\arabic{enumi}.}
\tightlist
\item
  A working workflow with explicit media evidence, decision ownership
  and correction rules.
\item
  Reproducible assembly and local-repair experiments, including frozen
  100-project validation and a retained direct-model comparison.
\item
  A narrated case showing where structural guarantees stop: correct
  assets still require explicit pause and caption conventions.
\end{enumerate}

These contributions establish technical support for correction. Their
effects on editing quality and creator effort remain to be evaluated.

\section{Related work}\label{related-work}

\textbf{Augmentation and mixed initiative.} REAL follows the
longstanding goal of amplifying human judgment
\citep{engelbart1962augmenting}. Mixed-initiative systems must decide
when to act and when to defer \citep{horvitz1999mixed}; human-centered
AI combines automation with human control \citep{shneiderman2020human}.
Evaluating an augmentation claim ultimately requires people working with
the system \citep{haupt2025centaur}. We test technical preservation here
and leave interaction benefits for a creator study.

\textbf{Interactive editing.} DemoCut structures instructional
recordings from lightweight marks \citep{chi2013democut}. QuickCut
connects annotations, transcripts and frame-level selection
\citep{truong2016quickcut}; B-Script recommends B-roll and insertion
positions within a transcript \citep{huber2019bscript}. ExpressEdit
supports natural-language and sketch references to temporal, spatial and
operational changes \citep{tilekbay2024express}. REAL likewise uses
language to address editing choices, emphasizing their persistence
across preparation, review, retiming and export.

\textbf{Agent workflows and readable artifacts.} LAVE combines footage
descriptions with agent assistance and direct manipulation
\citep{wang2024lave}. EditDuet uses Editor and Critic agents to
construct and assess sequences \citep{sandoval2025editduet}. Crayotter
exposes retrieval reports, analyses, blueprints and tool outputs in
long-form production \citep{yan2026crayotter}. These are close
precedents: staged agents and inspectable plans are not new
contributions here. REAL focuses on who owns a choice, which fields a
revision preserves, and independent checks of its translation into a
project file.

Its host uses transformer-based models \citep{vaswani2017attention}, but
REAL's representations are authored interfaces rather than learned
latent codes. Concept bottlenecks make intermediate concepts available
for intervention \citep{koh2020concept}. Music systems illustrate
several ways to keep a person in control: Piano Genie maps an
eight-button performance to piano music \citep{donahue2019piano}, while
Hookpad Aria offers continuations and infilling inside a lead-sheet
editor \citep{donahue2024aria}. YuE2 uses symbolic planning before audio
generation \citep{yuan2026yue2}. These interfaces and plans motivate
REAL's editable storyboard; their results do not establish benefits for
video editing. REAL applies this principle to media selection and
temporal placement without training a new model.

\textbf{Evaluation boundaries.} Message-driven editing benchmarks test
selection and sequence agreement \citep{ogata2026medit}, whereas
QVHighlights supplies annotated video moments
\citep{lei2021qvhighlights}. Such tasks differ from perceptual video
quality evaluation \citep{huang2024vbench}. Our controlled benchmark
isolates identity and timing against a team-authored specification. A
small retrieval probe checks a semantic component; neither substitutes
for evaluating people finding and repairing consequential mistakes.

\section{REAL: representations and
correction}\label{real-representations-and-correction}

REAL supports short narrated explanations assembled from footage,
diagrams and voice. Each stage produces a persistent artifact a creator
can inspect before translation continues (Figure 1). The design follows
documented production failures rather than a formal formative study:
expose a decision where its cause can be corrected, separate
observations from recommendations, and retain confirmed choices across
rebuilds.

\begin{figure}[htbp]
\centering
\includegraphics[width=1\linewidth,height=\textheight,keepaspectratio,alt={Four persistent artifacts separate narrative intent, media evidence, editorial decisions and finishing. Recorded word times retime placements while retaining confirmed choices.}]{architecture.pdf}
\caption{Four persistent artifacts separate narrative intent, media
evidence, editorial decisions and finishing. Recorded word times retime
placements while retaining confirmed choices.}
\end{figure}

A skill combines dialogue instructions, scripts and a handoff format.
The host selects the model; no foundation model is built into the four
skills. Models interpret media and requests; scripts parse the notebook,
probe files, bind assets, derive timing, apply revisions and validate
export. The structural tests use supplied descriptions and word times,
with no model calls.

\subsection{Intent and evidence}\label{intent-and-evidence}

A notebook groups typed components into \emph{beats}: narration spans
with associated visuals and sounds. Production notes stay outside the
narration clock. Requests can specify duration, asset tags and optional
status without project-file syntax. This prevents a stage direction from
silently becoming spoken time.

Cataloguing records duration, dimensions, frame rate, rotation and audio
streams. The production scanner samples 16 evenly spaced frame centres
by default; timestamped observations state their inspection coverage.
Model descriptions and suggested uses remain distinct from probed facts
and can be corrected before binding. Stable asset identifiers preserve
references after renaming.

The lexical binder combines explicit tags, declared associations and
weighted term overlap, enforcing media-kind compatibility and minimum
content evidence. Explicit tags receive priority; type compatibility
alone cannot establish a match. Contributing signals and alternatives
remain inspectable (Figure 2). Unresolved matches become visible
requests or placeholders. An optional grounding stage selects
source-event boundaries from ordered observations, records alternatives
and refuses unknown or uninspected evidence. Sampling does not establish
continuous coverage or precise motion boundaries.

Strongest matches bind first; notebook order and stable asset
identifiers break ties. These heuristics make a choice reproducible and
challengeable. Their semantic quality is not established by the tagged
assembly test.

\begin{figure}[htbp]
\centering
\includegraphics[width=1\linewidth,height=\textheight,keepaspectratio,alt={Five of 16 sampled frames from a Paganini lesson proxy and the logged path to an asset choice. Faces are masked for anonymous review. Probed facts, model observation and suggested use remain distinct. The tag and original excerpt window were supplied by the project script; the example does not demonstrate automatic passage discovery.}]{catalogue-path.pdf}
\caption{Five of 16 sampled frames from a Paganini lesson proxy and the
logged path to an asset choice. Faces are masked for anonymous review.
Probed facts, model observation and suggested use remain distinct. The
tag and original excerpt window were supplied by the project script; the
example does not demonstrate automatic passage discovery.}
\end{figure}

\subsection{Timing and decision
ownership}\label{timing-and-decision-ownership}

The storyboard is the authoritative edit. Each visual layer records
identity, asset, source interval, word anchor, requested duration, style
and selector. Audio lanes distinguish narration, source sound, music and
effects. The tracks view exposes selected assets and alternatives; edits
produce named revision commands rather than an unrecorded replacement of
the whole plan.

A \emph{word anchor} identifies a narration token whose time can change.
An estimated clock supports planning before recording; a measured clock
aligns words to a take. Updating the clock derives new placement times
without selecting content again. Creator ownership marks accepted
choices that \texttt{build\ -\/-preserve} retains; unrecognized
revisions and conflicting choices remain explicit errors.

For word \(i\), the default inline cue is \(s_i+0.4(e_i-s_i)+\delta\),
where \(s_i,e_i\) are word bounds and \(\delta\) is the creator's
offset. Base shots start at the beat boundary. Nearby cues may start
together within a 0.3-second window unless the creator marks one to keep
its separate timing; explicit holds extend a beat, while ordinary
overlays clip to its available span. The requested duration remains
separate so an earlier re-anchor can recover it. A pause policy can hold
the base shot until the next beat, subject to source length. These are
declared pacing conventions, not learned optimal choices.

Derived times retain precision until export. Absolute boundaries map to
frames with Python's ties-to-even rule; duration is the difference of
rounded end and start, not a separately rounded length. This avoids a
subtle defect: rounding 9.2498 seconds to milliseconds crosses a
half-frame boundary at 30 fps.

\subsection{Correction and export}\label{correction-and-export}

In the synthetic walkthrough, \texttt{B4:\ L3\ anchor\ =\ velvet} moves
the same \texttt{img07} layer from frames 607--631 to 562--586 at 30 fps
(Figure 3). Its requested duration remains 0.8 seconds; the other 15
non-gap placements keep their identities and boundaries, and the new
anchor survives rebuilding. The unusual words are test inputs. This
demonstrates the implemented correction path, not participant
performance.

\begin{figure}[htbp]
\centering
\includegraphics[width=1\linewidth,height=\textheight,keepaspectratio,alt={A stable placement across notebook, storyboard and FCPXML. The revision changes its anchor and ownership while retaining the asset. Labelled ranges are absolute; the exported offset is local to the containing clip.}]{placement-trace.pdf}
\caption{A stable placement across notebook, storyboard and FCPXML. The
revision changes its anchor and ownership while retaining the asset.
Labelled ranges are absolute; the exported offset is local to the
containing clip.}
\end{figure}

The compiler checks structure, assets, source bounds and text
classification before translating visuals, overlays, titles and audio.
Parity checks detect dropped elements; XML checks cover references,
frame alignment and timebases. It refuses unresolved decisions rather
than inventing editorial replacements. Target conventions still matter:
Final Cut positions are percentages of project height, requiring
conversion from pixel layout \citep{apple-animation}. A development
import exposed this error in both exporter and local reader, despite
their mutual agreement.

\subsection{Recorded-content case}\label{recorded-content-case}

A Paganini lesson complements the synthetic trace with spontaneous
speech. A review of the tracks page by the assisting model proposed
eight overlay re-anchors to associate captions and stills with named
phrases while retaining five base clips. These are proposed editorial
choices, not participant repairs. The final export contains the three
finished stills and supplied performance opening; its verification
script confirms that all 18 tutorial placement and source records shift
by exactly 419 frames for the intro. The creator reports that the
opening, stills, captions and audio play correctly in Final Cut,
including visible Glow. This is a single-project observation, not
systematic compatibility evidence.

The case exposed two upstream limits. A project adapter was needed to
express narration inside a supplied footage excerpt; downstream
source-window support does not by itself solve that notebook authoring
gap. PocketSphinx failed on two of five takes with out-of-vocabulary
words. On the other three, median onset disagreement with Whisper was
60--90 ms; the largest disagreement was 3.49 seconds on ``our'' in the
fourth take (7.83 versus 4.34 seconds). The project used Whisper word
times as its measured clock. Neither ASR source was checked against
manual acoustic annotations, so this is a documented fallback and
disagreement, not an aligner-accuracy benchmark.

\section{Evaluation}\label{evaluation}

We test whether declared choices survive compilation (RQ1), whether one
overlay can be corrected locally (RQ2), and which representation errors
remain on narrated content (RQ3). These establish technical
preconditions for creator intervention. Comprehension, usability and
perceived agency require people.

\subsection{Controlled assembly and local
repair}\label{controlled-assembly-and-local-repair}

A seeded generator produces four-to-eight-beat projects with tagged
assets, synthetic word times and zero-to-two overlays per beat. Sixteen
clips and twelve stills form the media pool. Tags remove semantic
ambiguity to isolate timing and compilation. The independent reference
imports no pipeline code but shares the team's declared semantics: a
beat spans its first to last word, base shots fill it from source frame
zero, overlays start 40\% into their anchor word and clip to the beat,
and absolute boundaries round to the 30 fps grid. It does not exercise
the separate hold-after-speech rule.

Initial development uses seeds 1--30; conformance checks use seeds
1001--1100. A changed-clock follow-up reuses the 100-project cohort and
exposes two further defects. After fixing them, we rerun assembly, the
scripted direct-model comparator and local repairs on one hash-frozen
pipeline. These are reused regression cohorts, not an untouched held-out
test. Measured word times are supplied before placement. Original-code
and estimated-clock runs check benchmark sensitivity and timing
dependence; they do not isolate individual implementation changes.

An independently implemented FCPXML reader recovers event identities and
boundaries. Greedy one-to-one matching requires the same asset or
normalized text and temporal intersection over union (IoU). We report
exact-placement F1 (IoU=1), averaged over projects, and the count of
all-exact projects. Failed compilation or parsing receives zero and
stays in the denominator. Exact matching requires the same identity and
both frame boundaries: a one-frame error leaves the prediction and
reference unmatched. Percentile 95\% intervals use 10,000 paired
project-seed bootstrap draws. These metrics exclude appearance and audio
mixing.

For each development seed, a fixed rule chooses one overlay and another
eligible anchor with enough duration remaining. Its revised bounds are
checked against an independent expectation. Collateral change counts
non-target event identities whose boundaries appear, disappear or
change. It does not measure effects, color, gain or human repair effort.

For the changed-clock check, one eligible overlay per seed is
re-anchored, then all supplied word times are scaled by 0.8 or 1.2 in
separate copies. Word identities and requested durations stay fixed. We
check complete reference placements and saved creator-choice fields
after retiming and a rebuild that preserves those choices, and compare
two compilations of each retimed story. This tests artificial timing
changes, not recordings of new speakers.

\subsection{Retained direct-model
comparison}\label{retained-direct-model-comparison}

Three recorded Claude CLI identifiers---Sonnet 5.5, Opus 5.5 and Fable
5.1---each have 30 retained attempts with and without additional FCPXML
format notes. Every prompt supplies the notebook, catalogue, words,
revisions and placement rules; models return XML with tools disabled and
no repair loop. Thus the task is model-as-compiler, not a workflow
without intermediate representations.

We audit all 180 prompts and response/XML pairs. Reference inputs and
catalogue asset entries match E10's revised development comparator.
Missing or unparseable XML receives zero temporal credit. Parsable XML
retains raw scores even when invalid; a secondary validity-gated score
assigns zero unless both Apple's 1.13 DTD and the existing semantic
validator pass. The reader covers all encountered constructs. Pointwise
paired intervals are exploratory, with no multiplicity adjustment or
equivalence margin. Regenerated projects and retained final attempts
make this a retrospective development comparison; CLI labels do not
establish immutable model revisions. Historical title-position
instructions also preclude an appearance comparison.

\subsection{Narrated reconstruction and semantic
probe}\label{narrated-reconstruction-and-semantic-probe}

A separate 30-second transformer explanation has an agent-prepared
reference cut: seven trimmed shots, seven caption images, ElevenLabs
narration and a frame-based edit list. REAL receives those labelled
assets, notebook and audio; PocketSphinx aligns the supplied words.
Saved E05 development modes use blank pauses (A) or hold shots through
pauses (B). The reference leads speech by 100 ms and places captions
from 400 ms after each cut to 200 ms before the next. Shot and caption
errors and missing picture use the full 900-frame horizon, including any
tail after the predicted timeline. Provider word times are a comparison
signal, not independently annotated acoustic truth.

A small QVHighlights probe \citep{lei2021qvhighlights} checks semantic
selection without reference tags. Five development and five
source-disjoint-at-selection queries use one frame per two seconds. REAL
describes frames without the query, then grounds the request in those
observations; direct vision receives images and the query. Both use the
recorded Sonnet 5.5 identifier, but REAL uses two calls versus one.
Official evaluator parity was checked across 1,550 sample predictions.
Fresh-case uncertainty bootstraps source videos with their queries.

\subsection{Reproduction boundaries}\label{reproduction-boundaries}

The artifact retains frozen implementations, input media, references,
predictions and per-case failures. The result table and numerical
summaries are generated from saved evidence. The live pipeline can
evolve independently of the evaluated snapshots; filenames and a
successful current build do not establish version parity. The one-shot
supplement includes actual prompts, responses and the executed
reader/scorer, with hashes checked before and after replay. An
independently extracted replay reproduces its summary exactly.

\section{Results}\label{results}

\subsection{Assembly and local repair}\label{assembly-and-local-repair}

The frozen measured-clock pipeline compiles and matches every one of
1,269 reference placements in 100/100 validation projects. This
establishes conformance to the declared generator. The benchmark exposed
five timing defects during development; the original implementation was
exact on 0/100 projects.

With estimated word times, 14 of 1,269 placements are exact (about 1\%);
project-averaged start/end errors are 0.31/0.66 seconds among the 641
placements matched at IoU≥.5 (maximum 2.37 seconds), with 628 unmatched
placements excluded from those error averages.

The imported implementation was exact on 28/30 development seeds.
Retaining derived precision resolved two one-frame discrepancies without
changing the reference or scorer. In the 30 development re-anchor
repairs, 30/30 targets have correct revised bounds and 0 non-target
identities change boundaries. This supports the tested locality
contract, not human repair efficiency.

Under ×0.8 and ×1.2 word-time scales, the correction survives retiming
and preserved rebuilding in 100/100 projects each, with saved choices
unchanged and repeated compilation byte-identical. The check exposed two
one-frame endpoint errors among 200 retimed exports and lost
style-ownership labels in 88 projects in the earlier code; both defects
were fixed before rerunning the reported snapshot.

\subsection{Direct generation from the same
plan}\label{direct-generation-from-the-same-plan}

Table 1 compares REAL with direct Claude CLI generation on all 30
development projects per condition. Raw temporal F1 retains parsable
invalid XML; gated F1 assigns zero unless both the DTD and semantic
validator pass. One retained response contains no XML and receives zero
in both columns.

\begin{table}[htbp]
\centering
\caption{REAL and direct Claude CLI generation from the same supplied plan (30 development projects per condition). Temporal F1 measures identity and boundaries; format pass requires both validators.}
\small
\begin{tabular}{@{}
  >{\raggedright\arraybackslash}p{(\linewidth - 10\tabcolsep) * \real{.22}}
  >{\raggedright\arraybackslash}p{(\linewidth - 10\tabcolsep) * \real{.16}}
  >{\raggedleft\arraybackslash}p{(\linewidth - 10\tabcolsep) * \real{.14}}
  >{\raggedleft\arraybackslash}p{(\linewidth - 10\tabcolsep) * \real{.16}}
  >{\raggedleft\arraybackslash}p{(\linewidth - 10\tabcolsep) * \real{.16}}
  >{\raggedleft\arraybackslash}p{(\linewidth - 10\tabcolsep) * \real{.16}}@{}}
\toprule\noalign{}
\begin{minipage}[b]{\linewidth}\raggedright
Compiler
\end{minipage} & \begin{minipage}[b]{\linewidth}\raggedright
Format notes
\end{minipage} & \begin{minipage}[b]{\linewidth}\raggedleft
All exact
\end{minipage} & \begin{minipage}[b]{\linewidth}\raggedleft
Format pass
\end{minipage} & \begin{minipage}[b]{\linewidth}\raggedleft
Temporal F1
\end{minipage} & \begin{minipage}[b]{\linewidth}\raggedleft
Gated F1
\end{minipage} \\
\midrule\noalign{}
REAL & Scripted & 30/30 & 30/30 & 1.000 & 1.000 \\
Sonnet 5.5 & Supplied & 27/30 & 30/30 & 0.991 & 0.991 \\
Opus 5.5 & Supplied & 30/30 & 30/30 & 1.000 & 1.000 \\
Fable 5.1 & Supplied & 30/30 & 30/30 & 1.000 & 1.000 \\
Sonnet 5.5 & Omitted & 24/30 & 22/30 & 0.948 & 0.726 \\
Opus 5.5 & Omitted & 26/30 & 26/30 & 0.964 & 0.867 \\
Fable 5.1 & Omitted & 28/30 & 26/30 & 0.958 & 0.867 \\
\bottomrule
\end{tabular}
\end{table}

With format notes, Opus and Fable match REAL on all 30 projects. The
REAL-minus-Sonnet difference is .009 {[}0.000, .020{]}. Direct
generation can therefore reproduce the supplied plan; equal sample
scores do not establish general equivalence. REAL makes validation and
rebuilding of saved decisions scripted operations requiring no new model
call. Without notes, paired differences are .052 {[}.016, .093{]}, .036
{[}.007, .072{]} and .042 {[}0.000, .117{]}. Format failures include
still-image element choices, frame-duration declarations and child
ordering. These are exploratory development comparisons. All 180
prompt/reference and response-extraction audits pass; a portable frozen
replay reproduces the summary byte-for-byte.

\subsection{Where conformance stops}\label{where-conformance-stops}

The narrated reconstruction selects all seven intended shots and caption
assets, but timing conventions differ. Blank-pause mode leaves 209 of
900 reference frames without a base picture and ends at frame 875.
Holding shots through pauses recovers all seven shot boundaries, while
mean caption start/end errors remain 8/6 frames. Source clips were
reference-trimmed: this does not demonstrate discovery of those
conventions in raw footage. Against provider word onsets, PocketSphinx
median absolute error is 51.3 ms (95th percentile 90.9 ms); neither
timestamp source is independent acoustic ground truth.

On the five fresh QVHighlights cases, REAL temporal and direct vision
each retrieve 3/5 at IoU .5, but mAP is 25.75 versus 32.83. The
exploratory source-bootstrap difference is -7.08 percentage points
{[}-24.33, 7.17{]}. The small, unequal-call probe establishes neither
superiority nor equivalence. Detailed predictions and official-evaluator
parity are in the artifact.

Frozen source hashes, all-case rows, independent references and
reproduction commands accompany the results. No new model inference was
used for this analysis. Native import and creator outcomes require their
own evidence.

\section{Discussion and limitations}\label{discussion-and-limitations}

\subsection{A contract for
intervention}\label{a-contract-for-intervention}

The experiments support a limited but useful claim: specified identities
and timing can survive compilation and local re-anchoring. Direct models
can also compile explicit plans accurately. REAL's proposed value is
therefore the addressable decision state and preservation contract, not
exclusive ability to produce correct XML. Whether creators understand
and benefit from that state is open.

Three design implications follow. \textbf{Expose the cause at an
editable boundary:} correct a media observation in the catalogue and a
placement in the storyboard. \textbf{State what a revision preserves:}
an anchor change should retain the asset and requested duration,
including who chose each saved value; the style-label defect shows that
ownership metadata can be lost even when values survive. Replacing an
asset may require new source coordinates. \textbf{Validate the claim
where it matters:} schema, temporal agreement and native appearance need
different evidence. The shared coordinate error shows why internal
agreement alone is insufficient.

\subsection{Limits of the evidence}\label{limits-of-the-evidence}

The team-authored generator covers tagged assets, synthetic words, 30
fps, zero source in-points and constrained overlays. New seeds sample
that grammar, not a population of creator projects. The reference shares
its specification with the implementation. One-shot outputs are retained
development attempts; detailed format notes encode expertise that the
compiler already contains. Equal sample scores establish neither general
equivalence nor unique need for scripts. Title appearance and native
import remain outside temporal F1. The repair and artificial-clock
checks cover re-anchoring on reused fixtures; preservation with new
recordings and other revision operations remains open.

The narrated case uses a synthetic voice and reference-trimmed assets.
Recovering its shot boundaries depends on the hold convention, while
captions remain mismatched. The real Paganini case needed an adapter and
alignment fallback; it is development experience, not a participant
study. Frame sampling can miss motion or sound, and plausible
descriptions can still be wrong. The retrieval probe is tiny,
single-model and unequal in call count; public videos may occur in
pretraining. We have no full-split QVHighlights result.

\subsection{Evaluation with creators}\label{evaluation-with-creators}

A counterbalanced study could compare correction through tracks and chat
using the same footage, model, initial errors and time budget
\citep{haupt2025centaur}. Tasks would cover incorrect observations,
source events and overlay times. Outcomes include localization, repair
success, time, collateral changes and perceived control. Readable state
could aid repair but impose a learning cost; neither expectation has
been tested.

\section{Conclusion}\label{conclusion}

REAL makes editorial choices inspectable and locally correctable. The
experiments show identity and timing preservation under declared
semantics; two direct models also compiled the supplied plans exactly.
Narrated reconstruction reveals the need to express pause and caption
conventions explicitly. Whether readable intermediate decisions help
creators locate and repair errors remains an empirical question.

\section*{Declaration on Generative
AI}\label{declaration-on-generative-ai}
\addcontentsline{toc}{section}{Declaration on Generative AI}

The authors directed the research ideas, experimental design and
scientific judgments, with human inspection of videos and Final Cut Pro
imports. ChatGPT/Codex, Claude and Gemini supported implementation
brainstorming, coding, debugging, and experiment planning, tracking and
coordination, alongside manuscript preparation, formatting,
paraphrasing, reference checking and citation management. ElevenLabs
supplied synthetic narration. The authors retain full responsibility for
the work.

\emergencystretch=3em
\bibliography{references}
\appendix
\emergencystretch=1em
\section{Implementation of the
skills}\label{implementation-of-the-skills}

The code is public at \texttt{github.com/This-Goober/REAL} under
\texttt{skills/batch-2/skills/}: four skills, 10,508 lines of Python,
standard library plus Pillow, ffmpeg and PocketSphinx. This appendix
states the representations and rules the results depend on, as
pseudo-code where a sentence would be ambiguous.

\subsection{Division of labour}\label{division-of-labour}

\begingroup\small\begin{longtable}[]{@{}
  >{\raggedright\arraybackslash}p{(\linewidth - 6\tabcolsep) * \real{0.1236}}
  >{\raggedright\arraybackslash}p{(\linewidth - 6\tabcolsep) * \real{0.2697}}
  >{\raggedright\arraybackslash}p{(\linewidth - 6\tabcolsep) * \real{0.3371}}
  >{\raggedright\arraybackslash}p{(\linewidth - 6\tabcolsep) * \real{0.2697}}@{}}
\caption{Who decides what. The model is whichever model the host runs;
the structural experiments supply descriptions and word times and call
none.}\tabularnewline
\toprule\noalign{}
\begin{minipage}[b]{\linewidth}\raggedright
Stage
\end{minipage} & \begin{minipage}[b]{\linewidth}\raggedright
The model
\end{minipage} & \begin{minipage}[b]{\linewidth}\raggedright
The scripts
\end{minipage} & \begin{minipage}[b]{\linewidth}\raggedright
Refuses when
\end{minipage} \\
\midrule\noalign{}
\endfirsthead
\toprule\noalign{}
\begin{minipage}[b]{\linewidth}\raggedright
Stage
\end{minipage} & \begin{minipage}[b]{\linewidth}\raggedright
The model
\end{minipage} & \begin{minipage}[b]{\linewidth}\raggedright
The scripts
\end{minipage} & \begin{minipage}[b]{\linewidth}\raggedright
Refuses when
\end{minipage} \\
\midrule\noalign{}
\endhead
\bottomrule\noalign{}
\endlastfoot
brainstorm & Drafts sequence and mock components with the creator & A
words-per-second ballpark, badged as such & Asked for media, ids or
times \\
catalogue & Describes sampled frames; suggests uses & Probes facts;
samples 16 frames per clip; groups near-duplicate takes by perceptual
hash; assigns stable ids & Sameness is uncertain (asks); unknown facts
stay \texttt{unresolved} \\
storyboarding & Interprets loose notes; reviews the tracks page; picks a
source event from listed observations & Parse, bind, clock, place,
revise, preserve, validate, render the tracks page & No notebook;
unknown revision; anchor text and index disagree; unclassified text; a
line that would change the clock \\
compile & Nothing & Adapt, parity, validate, align, assemble voice,
report & No reel; a dropped element; a validator error; any request for
a time \\
\end{longtable}\endgroup

Handoffs: \texttt{NOTEBOOK.md} to \texttt{notebook.json} to
\texttt{bindings.json} and \texttt{clock.json} to \texttt{reel.json} to
\texttt{story.json} to the timeline file. Each is validated against a
JSON schema on write and on read. The reel is the only editorial
handoff; the story is regenerated from it on every run and never edited.

\subsection{Vocabulary}\label{vocabulary}

{\def\LTcaptype{none} % do not increment counter
\begingroup\small\begin{longtable}[]{@{}
  >{\raggedright\arraybackslash}p{(\linewidth - 2\tabcolsep) * \real{0.1549}}
  >{\raggedright\arraybackslash}p{(\linewidth - 2\tabcolsep) * \real{0.8451}}@{}}
\toprule\noalign{}
\begin{minipage}[b]{\linewidth}\raggedright
Term
\end{minipage} & \begin{minipage}[b]{\linewidth}\raggedright
Meaning
\end{minipage} \\
\midrule\noalign{}
\endhead
\bottomrule\noalign{}
\endlastfoot
beat & One narration block with its visuals and sounds; spans first to
last word, extended by holds \\
component & One annotation such as
\texttt{\textless{}show\ image\ of\ two\ waves\textgreater{}}; typed
from its first four words; id is a hash of its text; modifiers
\texttt{\textasciitilde{}4s}, \texttt{!}, \texttt{?}, \texttt{\#tag};
\texttt{\textbackslash{}\textbar{}} means simultaneous \\
word anchor & Token index plus its text; placements store anchors, times
are derived; text and index must agree \\
clock & Word times, estimated from the syllable model or measured by
forced alignment; swapping re-derives every placement without
re-planning \\
layer & One visual on one lane: asset, source interval, anchor,
requested duration \texttt{dur}, occupied span, style,
\texttt{chosen\_by}, runner-ups \\
ownership & \texttt{chosen\_by:\ creator}, set by any revision;
preserved by id across rebuilds; a conflict is asked once \\
stub & A grey placeholder file at the exact size and duration for an
unbound component, so the import shows a hole, not missing media \\
\end{longtable}\endgroup
}

Revision lines, written by the tracks page on a click or typed by the
creator, one per instruction:

\begingroup\small\begin{verbatim}
B7: L1 anchor = pressure      B7: L1 source = 12-18      B7: L1 offset = +200
B7: layer L1 = A03            B7: remove L1              B7: confirm L1
B7: stagger L1                B7: add text = "..." anchor 150
B7: <free text>               -> a note on the layer or beat
B7: duration = 6.0 | hold longer | trim 2s   -> refused: the clock is not a placement
\end{verbatim}\endgroup

\subsection{Timing}\label{timing}

Estimated word duration is
\(\mathrm{speed}\,(W_0 + K_s\,\mathrm{syllables})\) with
\(W_0 = 0.0665\) s, \(K_s = 0.1403\) s, plus a pause after punctuation
(0.111 s comma, 0.148 s dash, 0.258 s sentence, 0.332 s paragraph, 0.406
s beat). Fitted to 486 recorded words in five sections: 3.11 words per
second, worst section error 5.8\%. Text durations (0.6 s plus 0.32 s per
word, minimum 1.2 s) are conventions, not measurements.

\begingroup\small\begin{verbatim}
derive(plan, words):                 # plan is fixed; only word times change
  for beat in plan:
    n0, n1 = start(first word), end(last word)
    beat.t0 = max(0, n0 - longest explicit cold open)
    leads: laid backwards from n0, one cursor per lane; flagged if they no longer fit
    holds: after n1, in written order
    inline p:    p.t0 = n0 if anchored on the first word, else cue(anchor) + offset
    post p:      p.t0 = n1 + offset
    beat.t1 = max(n1, end of holds, end of post layers)
    cued a, b consecutive by t0: if 0 < b.t0 - a.t0 < 0.3 and not b.stagger: b.t0 = a.t0
    base visual: runs to the next base or beat.t1, never under 0.4 s; explicit lengths fixed
    overlay:     t1 = min(t0 + requested, beat.t1), flagged clipped when cut
    every p:     span = t1 - t0; dur keeps the creator's explicit number
  if pause policy = hold: each base runs to 0.1 s before the next beat's first word
cue(i) = s_i + 0.4 (e_i - s_i) for word bounds s_i, e_i

swap(clock, measured):
  s_i, e_i = measured[i] where present (missing words keep estimates, warned)
  derive(plan, words)                # anchors untouched
  report per beat: estimated vs measured span; a coherent same-sign error is a clock fault
\end{verbatim}\endgroup

Measurement runs PocketSphinx forced alignment on the assembled take,
with a built-in out-of-vocabulary list for music terms and compounds
spelled from their parts; it fails rather than guess when audio and
script disagree.

\subsection{Frames}\label{frames}

\begingroup\small\begin{verbatim}
frames(t) = round_half_even(t * (den / num))   # rate formed first; called once per boundary
duration = frames(t1) - frames(t0)             # never a separately rounded length
\end{verbatim}\endgroup

Snapping happens once, in the adapter; the story carries integer frames
and both targets write exact rationals of the frame duration. Rounding
derived seconds to milliseconds first moved 9.2498 s across a half
frame; evaluating \texttt{t\ *\ 3000\ /\ 100} turned 8.55 s into
256.50000000000006, rounding to 257 against the reference's 256 (failure
14).

\subsection{Binding}\label{binding}

\begingroup\small\begin{verbatim}
score(c, a):                          # component c, catalogued asset a
  s = 0
  for tag in c.tags: s += 6 if tag in a.tags else 2.5 if tag in a.text else 0
  for field, cap in (binds_to 4, shows 3, role 2.5, use 2):
    s += min(cap, sum(idf(w) for shared w) / 2.2)        # idf(w) = log((N+1)/(df(w)+0.5))
  if c is text and c.content in a.excerpt: s += 3
  content = s                                            # evidence about content only
  s += kind(c.type, a.kind)   # +2 same; -1.5 / -2.5 image-video; -9 incompatible
  s -= 8 if a.reference, 3 if a.unused, 20 if a.do_not_use, 4 if file missing

bind(c): rank assets by s, ties by id
  text:      own content; an asset attaches only on a verbatim match
  narration: top if s >= 2 and content >= 1, confidence low (confirmed by ear)
  else:      top if s >= 2 and content >= 2.5; high iff s >= 5, margin >= 1.5, kinds agree
             otherwise a stub and a numbered request
\end{verbatim}\endgroup

Type words (image, clip, show, caption) are stop words. Kind alone never
binds. Runner-ups and signals are kept for the tracks page.

\subsection{Preservation}\label{preservation}

\begingroup\small\begin{verbatim}
rebuild(notebook, catalogue, old_reel):
  reel = assemble(bind(notebook, catalogue), derive(plan, words))
  for item in old_reel where chosen_by = creator:
    if item.asset exists: copy asset, anchor, dur, style, style_by, chosen_by into reel[item.id]
    else: ask once
retime(reel, words):
  creator-moved layers keep anchor and length; the rest take derive()'s times
  each beat: t1 = max(narration end, ends of holds and post layers, held pause)
apply(reel, lines):
  refuse a line that changes the clock; apply the rest; mark the item creator's; refit the beat
  report every unparsed line; exit non-zero if anything was refused or unhandled
\end{verbatim}\endgroup

A re-anchor restores \texttt{dur} when the beat has room, because
\texttt{dur} is kept apart from the span the layer last occupied
(failure 13). The same inputs compile to the same bytes.

\subsection{Compile}\label{compile}

\begingroup\small\begin{verbatim}
compile(reel, target):
  validate(reel): schema; anchor text equals the word at its index; a style on every
                  text layer; an asset or a stub on every other layer; audio kinds
  story = adapt(reel): snap boundaries once; resolve media facts; compute no time
  assert layers in = elements out and lanes in = clips out     # parity, or refuse
  file = target.build(story); target.validate(file) reports 0 errors or nothing is delivered
\end{verbatim}\endgroup

The FCPXML validator checks well-formedness and Apple's DTD, reference
and id integrity, every time on a whole frame, stills as \texttt{video}
with zero duration and timed media as \texttt{asset-clip}, a contiguous
spine, children inside the parent's timebase, in and out points inside
the source, and media files present. Children are emitted in the order
the DTD's own error text prescribes. Positions are converted from pixels
to percent of project height on both axes (failure 16).

\subsection{Failure modes the checks
catch}\label{failure-modes-the-checks-catch}

Table 4 lists the failure modes the conformance checks are built to
catch and the rule that prevents each. Every one produces a
valid-looking timeline or is caught only by the compile step's refusal.

\begingroup\small\begin{longtable}[]{@{}
  >{\raggedright\arraybackslash}p{(\linewidth - 8\tabcolsep) * \real{0.0220}}
  >{\raggedright\arraybackslash}p{(\linewidth - 8\tabcolsep) * \real{0.2857}}
  >{\raggedright\arraybackslash}p{(\linewidth - 8\tabcolsep) * \real{0.1209}}
  >{\raggedright\arraybackslash}p{(\linewidth - 8\tabcolsep) * \real{0.2637}}
  >{\raggedright\arraybackslash}p{(\linewidth - 8\tabcolsep) * \real{0.3077}}@{}}
\caption{Failure modes and the rule that prevents each. The assembly
benchmark catches 1 to 5 and 13, the narrated reconstruction 7 to 12,
the retiming check 14 and 15, the shipped example 6, and an editor
import 16.}\tabularnewline
\toprule\noalign{}
\begin{minipage}[b]{\linewidth}\raggedright
\#
\end{minipage} & \begin{minipage}[b]{\linewidth}\raggedright
Failure
\end{minipage} & \begin{minipage}[b]{\linewidth}\raggedright
Kind
\end{minipage} & \begin{minipage}[b]{\linewidth}\raggedright
Symptom
\end{minipage} & \begin{minipage}[b]{\linewidth}\raggedright
Rule
\end{minipage} \\
\midrule\noalign{}
\endfirsthead
\toprule\noalign{}
\begin{minipage}[b]{\linewidth}\raggedright
\#
\end{minipage} & \begin{minipage}[b]{\linewidth}\raggedright
Failure
\end{minipage} & \begin{minipage}[b]{\linewidth}\raggedright
Kind
\end{minipage} & \begin{minipage}[b]{\linewidth}\raggedright
Symptom
\end{minipage} & \begin{minipage}[b]{\linewidth}\raggedright
Rule
\end{minipage} \\
\midrule\noalign{}
\endhead
\bottomrule\noalign{}
\endlastfoot
1 & Base shot cued 40\% into its first word & placement & Every beat
opens 3 to 4 frames late & A cue on a line's first word lands on the
beat boundary \\
2 & A timed still after a running shot taken as the next base & notebook
semantics & Shot cut short, up to a third of the beat black & An image
interrupting a running shot is an overlay, not a cut \\
3 & Explicit \texttt{\textasciitilde{}1.7s} overwritten by the span the
first clock allows & state & Clipped durations stay clipped after
re-timing & \texttt{dur} is the creator's number; the span is recomputed
per clock \\
4 & Instructions say ``measure first''; the code can only estimate, then
swap & instruction-code drift & Inline versus hold decided on estimated
lengths & The clock takes measured times directly; every stated rule has
a regression check \\
5 & Re-anchoring a hold leaves the beat stretched & state & Beat
overlaps the next; compile refuses & The beat end is recomputed from its
contents after every revision \\
6 & Validation splits narration on whitespace, the clock with a regex &
interface & A valid notebook fails validation & One shared tokeniser \\
7 & A visual description ending in ``.'' typed as narration & type
inference & Descriptions counted as speech, 2.0 s on a 60 s example &
Type from the first four words; a component is never speech \\
8 & The aligner writes \texttt{s}/\texttt{e}; the clock reads
\texttt{start}/\texttt{end} & interface & The measured path cannot
consume the aligner's output & A schema on the measured handoff; the
reader accepts both spellings \\
9 & Technical compounds absent from the aligner's dictionary & algorithm
& Alignment refuses technical vocabulary & Compounds pronounced from
their parts; a built-in out-of-vocabulary list \\
10 & Prose before the first heading becomes an empty beat &
representation & Zero-length beat; compile refuses & Leading prose is a
note, not a beat \\
11 & Silence between lines is always a black gap & representation & 209
of 900 frames black against an editor's cut & The \texttt{hold} pause
policy \\
12 & A held shot outlasts its clip & representation & A one-frame
alignment difference makes compile refuse & The cut moves to where the
footage ends, flagged; a clip shorter than its line is refused \\
13 & Re-anchoring a clipped cue keeps the clipped length & state &
Requested duration lost after a move & \texttt{dur} kept apart from the
occupied span \\
14 & Half-frame rounding, 8.55 s to 256.50000000000006 frames &
arithmetic & One-frame errors on exact half frames & Rate formed first;
one snap per boundary \\
15 & Style ownership relabelled as model-selected on rebuild & state &
Values survive, provenance lost & \texttt{style\_by} on the preservation
list \\
16 & Positions written in pixels; Final Cut expects percent of height &
target & Overlays off-position while exporter and reader agree &
Conversion at the target \\
\end{longtable}\endgroup

Three rules generalise the table. A skill is two artifacts, instructions
and scripts, so every rule stated in an instruction file has a
regression check. Handoffs drift without types, so every handoff has a
schema and the skills share one core module. Where a fixed rule would
stand in for something the format cannot say, the notebook gets the word
for it, and what remains is a declared convention.

\subsection{Tests}\label{tests}

A 40-check end-to-end suite runs the scripts on synthetic data, one
check per stated rule (35 rules), and 38 unit tests with 48 subtests
cover timing precision, layout, geometry, titles, style ownership and
grounding. All pass.

\subsection{Second target}\label{second-target}

The sequencing logic never learns an output format: the adapter produces
the story, and a target is a module with \texttt{build} and
\texttt{validate}. Two targets exist. FCPXML is the primary.
OpenTimelineIO, which DaVinci Resolve and Premiere Pro import natively,
is mapped element for element (beats to the first video track, overlay
lanes to further tracks, four audio lanes to four tracks). It has no
title schema, so text travels as a generator clip with its parameters
or, for editors without a title importer, as a transparent PNG drawn at
sequence size and calibrated to Final Cut's template (1.78 em per
nominal em in a 1080 by 1920 project, block 0.2 em below the glyph
centre; 1.73 measured on an export). All 130 benchmark projects read
back equivalent through the two independent readers. The Paganini
tutorial imported into both editors renders to 2,747 frames in each, all
6 cuts and 22 caption edges on the expected frame, glyph boxes within
1\%, median SSIM 0.992 between renders, audio lag 0 ms; a
semi-transparent still is 22\% darker in Resolve, a renderer property.
One tutorial in two editors, not systematic compatibility evidence.

\end{document}
